\section{Conceptos y Teorías}

    \textbf{Teoría de Brønsted-Lowry:} Un ácido es una especie que dona protones ($H^{+}$) y una base es una especie que los acepta.


\subsection{Fuerza de ácidos y bases}
\begin{itemize}
    \item \textbf{Ácidos fuertes:} Se disocian completamente en agua.
    \item \textbf{Ácidos débiles:} Se disocian parcialmente.
    \item \textbf{Bases fuertes:} Se disocian completamente.
    \item \textbf{Bases débiles:} Se disocian parcialmente.
\end{itemize}

\subsection{Ionización y equilibrio}
El fenómeno de ionización de ácidos y bases débiles se describe mediante una constante de equilibrio (\textit{K}) que indica la extensión de la reacción.

\section{Constante de concentración ($K$)}

\begin{equation*}
    K=\frac {[C]^c*[D]^d}{[A]^a*[B]^b}
\end{equation*}

Los sólidos y líquidos no cambian en su [x] inicial y final, por lo tanto sus valores son 1 y por ello no se escriben en la ecuación de equilibrio.

\section{Ácidos y bases según Brønsted-Lowry}

\begin{itemize}
    \item Ácidos: Donan protones $[H^{+}]$
    \item Bases: Aceptan protones $[H^{+}]$ 
\end{itemize}

Se denominan ácidos fuertes a aquellos ácidos que se disocian de manera completa. Estos son: 

\begin{itemize}
    \item $HCl$
    \item $HNO_{3}$
    \item $H_{2}SO_{4}$
    \item $HBr$
    \item $HI$
    \item $HClO_{4}$
\end{itemize}

El resto de ácidos se denominan ácidos débiles.

\section{Valores de $K_w$ y pH}

\begin{align*}
    K_w &= [H_{3}O^{+}]*[OH^{-}] \\
    K_w &= 1 \times 10^{-14} \\
    pH &= -\log[H_{3}O^{+}] \\
    pOH &= -\log[OH^{-}] \\
    pH + pOH &= 14
\end{align*}

\resumebox{ejercicio}{Cálculo de pH}{\footnotesize
    Calcula el pH de una solución de HCl 0.01 M y de una solución de NaOH 0.001 M. Determina también la concentración de $[OH^{-}]$ para la primera y de $[H_{3}O^{+}]$ para la segunda.
}

\resumebox{dato}{pH en alimentos}{\footnotesize
    El pH es un factor crítico en la conservación de alimentos, ya que afecta el crecimiento microbiano. Por ejemplo, la mayoría de las bacterias patógenas no crecen a pH menores de 4.6, por lo que los alimentos ácidos (como los enlatados) son más seguros.
}